% !TEX root = ../main.tex
\section{Proofs for Appendix~\ref{sec:upper-bound}}
\label{app:curvature-smoothness}

\begin{proof}[Proof of the curvature part of Proposition~\ref{prop:curvature-smoothness}]
For fixed $\blmbd$, the reduced payoff satisfies
\[
    \nabla_{\bx}\mathcal D_\rho(\bx,\blmbd)
    =
    \nabla f(\bx)+J_{\calA}(\bx)^\top\bq(\bx,\blmbd),
\]
where
\[
    \bq(\bx,\blmbd)
    =
    \operatorname{prox}_{\rho h^*}
    \bigl(\blmbd+\rho\calA(\bx)\bigr).
\]

The nonexpansiveness of the proximal mapping, the bounded Jacobian, and
the Lipschitz Jacobian imply
\[
    \|\bq(\bx,\blmbd)-\bq(\bx',\blmbd)\|
    \le \rho G_{\calA}\|\bx-\bx'\|.
\]
Assumption~\ref{ass:h} gives
\(\|\bq(\bx,\blmbd)\|\le L_h\).  Therefore
\[
    \|\nabla_{\bx}\mathcal D_\rho(\bx,\blmbd)
    -\nabla_{\bx}\mathcal D_\rho(\bx',\blmbd)\|
    \le L_{\calD,\bx}\|\bx-\bx'\|.
\]

Thus $\mathcal D_\rho(\cdot,\blmbd)$ is
$L_{\calD,\bx}$-weakly convex.  Adding
$\frac{1}{2\gamma}\|\bx-\bz\|^2$ gives
$\mu_\gamma$-strong convexity in $\bx$.

For fixed $(\bx,\bz,\bmu)$, $\mathcal D_\rho(\bx,\cdot)$ is concave in
$\blmbd$ because it is the pointwise infimum of affine functions of
$\blmbd$.  The term
$-\frac{1}{2s}\|\blmbd-\bmu\|^2$ is $1/s$-strongly concave.  Hence
$\mathcal P_s$ is $1/s$-strongly concave in $\blmbd$.
\end{proof}

\begin{proof}[Proof of the smoothness part of Proposition~\ref{prop:curvature-smoothness}]
The $\bx$-block bound follows from the Lipschitz estimate for
$\nabla_{\bx}\mathcal D_\rho$ in the proof of
Proposition~\ref{prop:curvature-smoothness} and the quadratic term
$\frac{1}{2\gamma}\|\bx-\bz\|^2$.

For the $\blmbd$-block, the identity
\[
    \nabla_{\blmbd}\mathcal P_s(\bx,\blmbd;\bz,\bmu)
    =
    \frac{1}{\rho}
    \bigl(\bq(\bx,\blmbd)-\blmbd\bigr)
    -\frac{1}{s}(\blmbd-\bmu)
\]
and the nonexpansiveness of
$\operatorname{prox}_{\rho h^*}$ give
$L_\lambda(s)=1/\rho+1/s$.
The same estimate, together with
\[
    \|\calA(\bx)-\calA(\bx')\|
    \le G_{\calA}\|\bx-\bx'\|,
\]
gives both cross Lipschitz bounds with constant $G_{\calA}$.
\end{proof}

\begin{proof}[Proof of Lemma~\ref{lem:dual-value-gain}]
Set
\[
    g^k
    :=
    \nabla_{\blmbd}\mathcal P_k
    (\bx^{k+1},\blmbd^k,\bz^k),
    \qquad
    \Delta_{\blmbd}^{k+1}=\tau_k g^k,
\]
where the second identity follows from the dual update.
By the $L_{d,k}$-smoothness of $d_k(\cdot;\bz^k)$,
\[
\begin{aligned}
&d_k(\blmbd^{k+1};\bz^k)-d_k(\blmbd^k;\bz^k) \\
&\quad\ge
\left\langle
\nabla_{\blmbd}d_k(\blmbd^k;\bz^k),
\Delta_{\blmbd}^{k+1}
\right\rangle
-\frac{L_{d,k}}{2}
\|\Delta_{\blmbd}^{k+1}\|^2.
\end{aligned}
\]
Danskin's theorem and Lemma~\ref{lem:v2-value-regularity} give
\[
\nabla_{\blmbd}d_k(\blmbd^k;\bz^k)
=
\nabla_{\blmbd}\mathcal P_k
\bigl(\bx_k(\blmbd^k;\bz^k),\blmbd^k,\bz^k\bigr).
\]
The cross-Lipschitz estimate in
Proposition~\ref{prop:curvature-smoothness} therefore yields
\[
\begin{aligned}
&\left\langle
\nabla_{\blmbd}d_k(\blmbd^k;\bz^k),
\Delta_{\blmbd}^{k+1}
\right\rangle \\
&\quad=
\tau_k\left\langle
\nabla_{\blmbd}\mathcal P_k
\bigl(\bx_k(\blmbd^k;\bz^k),\blmbd^k,\bz^k\bigr),g^k
\right\rangle \\
&\quad\ge
\tau_k\|g^k\|^2
-\tau_kG_{\cal A}
\|\bx^{k+1}-\bx_k(\blmbd^k;\bz^k)\|\|g^k\|.
\end{aligned}
\]
Using
\[
    \tau_k\|g^k\|^2
    =\frac{1}{\tau_k}\|\Delta_{\blmbd}^{k+1}\|^2,
    \qquad
    \tau_k\|g^k\|
    =\|\Delta_{\blmbd}^{k+1}\|,
\]
we obtain
\begin{equation}
\begin{aligned}
&d_k(\blmbd^{k+1};\bz^k)-d_k(\blmbd^k;\bz^k) \\
&\quad\ge
\left(\frac{1}{\tau_k}-\frac{L_{d,k}}{2}\right)
\|\Delta_{\blmbd}^{k+1}\|^2
-G_{\cal A}\|\bx^{k+1}-\bx_k(\blmbd^k;\bz^k)\|
\|\Delta_{\blmbd}^{k+1}\|.
\end{aligned}
\label{eq:dual-value-gain-appendix}
\end{equation}

Let
\[
    b_k:=
    \frac{1}{\tau_k}+\frac{1}{2s_k}-L_{d,k}.
\]
Multiplying \eqref{eq:dual-value-gain-appendix} by $2$ and rearranging gives
\[
\begin{aligned}
&\left(\frac{1}{\tau_k}-\frac{1}{2s_k}\right)
\|\Delta_{\blmbd}^{k+1}\|^2 \\
&\quad\le
2\bigl[d_k(\blmbd^{k+1};\bz^k)-d_k(\blmbd^k;\bz^k)\bigr]
-b_k\|\Delta_{\blmbd}^{k+1}\|^2 \\
&\qquad+
2G_{\cal A}\|\bx^{k+1}-\bx_k(\blmbd^k;\bz^k)\|
\|\Delta_{\blmbd}^{k+1}\|.
\end{aligned}
\]
The assumed parameter condition implies $b_k\ge L_{d,k}>0$. Young's inequality
with parameter $L_{d,k}$ gives
\[
2G_{\cal A}\|\bx^{k+1}-\bx_k(\blmbd^k;\bz^k)\|
\|\Delta_{\blmbd}^{k+1}\|
\le
\frac{L_{d,k}}{2}\|\Delta_{\blmbd}^{k+1}\|^2
+\frac{2G_{\cal A}^2}{L_{d,k}}
\|\bx^{k+1}-\bx_k(\blmbd^k;\bz^k)\|^2.
\]
Since $b_k\ge L_{d,k}$, we have
\[
    -b_k\|\Delta_{\blmbd}^{k+1}\|^2
    +\frac{L_{d,k}}{2}\|\Delta_{\blmbd}^{k+1}\|^2
    \le
    -\frac{L_{d,k}}{2}\|\Delta_{\blmbd}^{k+1}\|^2.
\]
Finally, the $\mu_\gamma$-strong convexity of
$\mathcal P_k(\cdot,\blmbd^k;\bz^k)$ gives
\[
\|\bx^{k+1}-\bx_k(\blmbd^k;\bz^k)\|^2
\le
\frac{2}{\mu_\gamma}
\left[
\mathcal P_k(\bx^{k+1},\blmbd^k;\bz^k)
-d_k(\blmbd^k;\bz^k)
\right].
\]
Substituting this estimate proves the first inequality.
\end{proof}

\begin{proof}[Proof of the second estimate in Lemma~\ref{lem:dual-value-gain}]
Fix $k\ge0$. By Proposition~\ref{prop:curvature-smoothness},
$\mathcal P_k(\cdot,\blmbd^k;\bz^k)$ is
$\mu_\gamma$-strongly convex. Applying the strong-convexity inequality
at $\bx^{k+1}$ to the minimizer $\bx_k(\blmbd^k;\bz^k)$ gives
\[
\begin{aligned}
&\mathcal P_k(\bx^{k+1},\blmbd^k;\bz^k)-d_k(\blmbd^k;\bz^k) \\
&\quad\le
\left\langle
\nabla_{\bx}\mathcal P_k(\bx^{k+1},\blmbd^k;\bz^k),
\bx^{k+1}-\bx_k(\blmbd^k;\bz^k)
\right\rangle \\
&\qquad
-\frac{\mu_\gamma}{2}
\|\bx^{k+1}-\bx_k(\blmbd^k;\bz^k)\|^2 \\
&\quad\le
\frac{1}{2\mu_\gamma}
\|\nabla_{\bx}\mathcal P_k(\bx^{k+1},\blmbd^k;\bz^k)\|^2,
\end{aligned}
\]
where the last inequality follows from Young's inequality.
The $\bx$-update and the $L_x$-Lipschitz continuity of
$\nabla_{\bx}\mathcal P_k(\cdot,\blmbd^k;\bz^k)$ imply
\[
\begin{aligned}
&\|\nabla_{\bx}\mathcal P_k(\bx^{k+1},\blmbd^k;\bz^k)\| \\
&\quad\le
\|\nabla_{\bx}\mathcal P_k(\bx^k,\blmbd^k;\bz^k)\|
+L_x\|\Delta_{\bx}^{k+1}\| \\
&\quad=
\left(\frac{1}{\eta_k}+L_x\right)
\|\Delta_{\bx}^{k+1}\|.
\end{aligned}
\]
Combining these estimates yields
\begin{equation}
\label{eq:primal-gap-control}
\mathcal P_k(\bx^{k+1},\blmbd^k;\bz^k)
-d_k(\blmbd^k;\bz^k)
\le
\frac{(1/\eta_k+L_x)^2}{2\mu_\gamma}
\|\Delta_{\bx}^{k+1}\|^2.
\end{equation}
Substituting \eqref{eq:primal-gap-control} into the first inequality of
Lemma~\ref{lem:dual-value-gain} gives a bound on the dual ascent term,
where the coefficient of
$\|\Delta_{\bx}^{k+1}\|^2$ becomes
\[
\frac{4G_{\mathcal A}^2}{L_{d,k}\mu_\gamma}
\frac{(1/\eta_k+L_x)^2}{2\mu_\gamma}
=
\frac{2G_{\mathcal A}^2(1/\eta_k+L_x)^2}
{L_{d,k}\mu_\gamma^2}.
\]
Substitute the resulting bound on
$(1/\tau_k-1/(2s_k))\|\Delta_{\blmbd}^{k+1}\|^2$
into \eqref{eq:basic-one-step} and combine the coefficients of
$\|\Delta_{\bx}^{k+1}\|^2$. The $\bz$-step term and both parameter-drift
terms remain unchanged. This proves \eqref{eq:dual-value-full-step}.
\end{proof}

\begin{proof}[Proof of Lemma~\ref{lem:partial-lyapunov-z-step}]
Subtract
$2[d_k(\blmbd^{k+1};\bz^k)-d_k(\blmbd^k;\bz^k)]$
from both sides of \eqref{eq:dual-value-full-step}. To replace
$d_k(\blmbd^{k+1};\bz^k)$ on the left by
$d_k(\blmbd^{k+1};\bz^{k+1})$, add
$-2[d_k(\blmbd^{k+1};\bz^{k+1})-d_k(\blmbd^{k+1};\bz^k)]$
to both sides.
By the minimum-value definition,
\[
d_k(\blmbd^{k+1};\bz^k)
\le
\mathcal P_k\bigl(\bx_k(\blmbd^{k+1};\bz^{k+1}),
\blmbd^{k+1};\bz^k\bigr).
\]
Subtracting this inequality from the corresponding equality at
$\bz^{k+1}$ and expanding the quadratic terms yields
\[
\begin{aligned}
&d_k(\blmbd^{k+1};\bz^{k+1})-d_k(\blmbd^{k+1};\bz^k) \\
&\quad\ge
\frac1\gamma\left\langle
\bz^k-\bx_k(\blmbd^{k+1};\bz^{k+1}),
\Delta_{\bz}^{k+1}\right\rangle
+\frac1{2\gamma}\|\Delta_{\bz}^{k+1}\|^2 \\
&\quad=
-\frac{\alpha_k(2-\alpha_k)}{2\gamma}
\|\bx^{k+1}-\bz^k\|^2 \\
&\qquad+
\frac{\alpha_k}{\gamma}\left\langle
\bx^{k+1}-\bx_k(\blmbd^{k+1};\bz^{k+1}),
\bx^{k+1}-\bz^k\right\rangle.
\end{aligned}
\]
Multiply this inequality by $-2$ and substitute it into the preceding
estimate. Its squared-step term combines with the $\bz$-step term in
\eqref{eq:dual-value-full-step} to give
$\alpha_k(2-\alpha_k)\|\bx^{k+1}-\bz^k\|^2/(2\gamma)$.
Both parameter-drift terms remain unchanged, proving
\eqref{eq:partial-lyapunov-z-step}.
\end{proof}

\begin{proof}[Proof of Lemma~\ref{lem:z-step-p}]
By Lemma~\ref{lem:v2-value-regularity}, $p_k$ has a
$1/\gamma$-Lipschitz continuous gradient and
\[
    \nabla p_k(\bz^k)
    =
    \frac1\gamma\bigl(\bz^k-\bx_k^\star(\bz^k)\bigr).
\]
The $\bz$-update gives
\[
    \bz^{k+1}-\bz^k
    =\alpha_k(\bx^{k+1}-\bz^k).
\]
Applying the smoothness inequality for $p_k$ and expanding the inner
product yields \eqref{eq:z-step-p}.
\end{proof}

\begin{proof}[Proof of Lemma~\ref{lem:augmented-one-step}]
By Lemma~\ref{lem:z-step-p},
\begin{equation}
\begin{aligned}
2\bigl[p_k(\bz^{k+1})-p_k(\bz^k)\bigr]
\le{}&
-\frac{\alpha_k(2-\alpha_k)}{\gamma}
\|\bx^{k+1}-\bz^k\|^2 \\
&+
\frac{2\alpha_k}{\gamma}
\left\langle
\bx^{k+1}-\bx_k^\star(\bz^k),
\bx^{k+1}-\bz^k
\right\rangle.
\end{aligned}
\end{equation}
Adding this inequality to
\eqref{eq:partial-lyapunov-z-step}, the quadratic terms combine as
\[
\frac{\alpha_k(2-\alpha_k)}{2\gamma}
-\frac{\alpha_k(2-\alpha_k)}{\gamma}
=
-\frac{\alpha_k(2-\alpha_k)}{2\gamma}.
\]
The two inner products combine into
\[
\frac{2\alpha_k}{\gamma}
\left\langle
\bx_k(\blmbd^{k+1};\bz^{k+1})
-\bx_k^\star(\bz^k),
\bx^{k+1}-\bz^k
\right\rangle.
\]
Substitution into \eqref{eq:partial-lyapunov-z-step} proves
\eqref{eq:augmented-one-step}.
\end{proof}

\begin{proof}[Proof of Lemma~\ref{lem:current-multiplier-boundedness}]
The proximal optimality condition gives
$\bq(\bx,\blmbd)\in\partial h(\by(\bx,\blmbd))$, and hence
$\|\bq(\bx,\blmbd)\|\le L_h$. Substituting
$\calA(\bx)-\by(\bx,\blmbd)
=(\bq(\bx,\blmbd)-\blmbd)/\rho$ into the dual update yields
\[
\blmbd^{k+1}
=\left[1-\tau_k\left(\frac1\rho+\frac1{s_k}\right)\right]\blmbd^k
+\frac{\tau_k}{\rho}\bq(\bx^{k+1},\blmbd^k)
+\frac{\tau_k}{s_k}\bmu^k.
\]
The coefficients are nonnegative and sum to one. Together with
$\bmu^{k+1}=(1-\beta_k)\bmu^k+\beta_k\blmbd^{k+1}$, this proves
$\|\blmbd^k\|,\|\bmu^k\|\le M$ by induction from the definition of $M$.
The center update then gives
$\|\bmu^{k+1}-\bmu^k\|\le2M\beta_k$.

At the saddle point, the dual first-order condition gives
\[
\blmbd_k^\star(\bz)
=\frac{s_k}{\rho+s_k}
  \bq\bigl(\bx_k^\star(\bz),\blmbd_k^\star(\bz)\bigr)
+\frac{\rho}{\rho+s_k}\bmu^k.
\]
Since both vectors in this convex combination have norm at most $M$,
$\|\blmbd_k^\star(\bz)\|\le M$.
\end{proof}

\begin{proof}[Proof of Lemma~\ref{lem:saddle-center-stability}]
Fix $k$ and $\bz$. Strong convexity gives the unique minimizer
$\bx_k(\blmbd;\bz)$ for each $\blmbd$. The finite, smooth,
$1/s_k$-strongly concave function $d_k(\cdot;\bz)$ tends to $-\infty$
as $\|\blmbd\|\to\infty$, so it has a unique maximizer.
The minimax equality and strong convexity in $\bx$ identify the unique
saddle pair and give
$\bx_k^\star(\bz)=\bx_k(\blmbd_k^\star(\bz);\bz)$.

The maximizing property at $\bx_k^\star(\bz)$ and strong convexity at
fixed $\blmbd$ imply
\[
\begin{aligned}
p_k(\bz)-d_k(\blmbd;\bz)
&\ge
\mathcal P_k(\bx_k^\star(\bz),\blmbd;\bz)
-\mathcal P_k(\bx_k(\blmbd;\bz),\blmbd;\bz) \\
&\ge\frac{\mu_\gamma}{2}
\|\bx_k(\blmbd;\bz)-\bx_k^\star(\bz)\|^2.
\end{aligned}
\]
This proves \eqref{eq:current-saddle-stability}.
For the center sensitivity of $\bx_k(\blmbd;\bz)$, add the
strong-convexity inequalities at the minimizers for $\bz$ and $\bz'$.
The only center-dependent term is $\|\bx-\bz\|^2/(2\gamma)$, so
\[
\mu_\gamma\|\bx_k(\blmbd;\bz')-\bx_k(\blmbd;\bz)\|^2
\le\frac1\gamma\left\langle
\bz'-\bz,\bx_k(\blmbd;\bz')-\bx_k(\blmbd;\bz)\right\rangle.
\]
For the saddle primal points, use the same argument for
$\max_{\blmbd}\mathcal P_k(\cdot,\blmbd;\bz)$, whose strong-convexity
modulus is also $\mu_\gamma$. Cauchy--Schwarz gives
\eqref{eq:current-center-stability} in both cases.
The latter argument also gives
\[
\left\langle\bz'-\bz,\bx_k^\star(\bz')-\bx_k^\star(\bz)\right\rangle
\ge\gamma\mu_\gamma
\|\bx_k^\star(\bz')-\bx_k^\star(\bz)\|^2.
\]
Together with the Danskin formula for $\nabla p_k$ and
$2\gamma\mu_\gamma>1$, this verifies the smoothness bound used above:
\[
\begin{aligned}
\gamma^2\|\nabla p_k(\bz')-\nabla p_k(\bz)\|^2
&\le\|\bz'-\bz\|^2
-(2\gamma\mu_\gamma-1)
\|\bx_k^\star(\bz')-\bx_k^\star(\bz)\|^2 \\
&\le\|\bz'-\bz\|^2.
\end{aligned}
\]
\end{proof}

\begin{proof}[Proof of Lemma~\ref{lem:moving-three-value-descent}]
Let $u_{k+1}:=(\bx^{k+1},\blmbd^{k+1},\bz^{k+1})$ and define
\[
\widehat\Psi_{k+1|k}:=
\mathcal P_{k+1}(\bx^{k+1},\blmbd^{k+1};\bz^{k+1})
-2d_k(\blmbd^{k+1};\bz^{k+1})+2p_k(\bz^{k+1}).
\]
Then
\[
\Psi_{k+1}(u_{k+1})-\Psi_k(u_k)
=\widehat\Psi_{k+1|k}-\Psi_k(u_k)+R_k,
\]
where
\[
R_k:=2\bigl[d_k(\blmbd^{k+1};\bz^{k+1})
-d_{k+1}(\blmbd^{k+1};\bz^{k+1})\bigr]
+2\bigl[p_{k+1}(\bz^{k+1})-p_k(\bz^{k+1})\bigr].
\]

Apply Lemma~\ref{lem:augmented-one-step} to the first difference.
Decompose the indefinite term at $\bx_k^\star(\bz^{k+1})$ and apply
Lemma~\ref{lem:saddle-center-stability}, followed by Young's inequality.
The four bounds in \eqref{eq:parameter-choice} absorb the resulting
cross terms and leave
\[
\begin{aligned}
\widehat\Psi_{k+1|k}-\Psi_k(u_k)
\le{}&
-\frac1{4\eta_k}\|\Delta_{\bx}^{k+1}\|^2
-\frac1{8\tau_k}\|\Delta_{\blmbd}^{k+1}\|^2
-\frac{\alpha_k}{2\gamma}\|\bx^{k+1}-\bz^k\|^2 \\
&+\frac{\beta_k(2-\beta_k)}{2s_k}
\|\blmbd^{k+1}-\bmu^k\|^2
\\
&+\frac12\left(\frac1{s_k}-\frac1{s_{k+1}}\right)
\|\blmbd^{k+1}-\bmu^{k+1}\|^2.
\end{aligned}
\]

It remains to bound $R_k$. Lemma~\ref{lem:current-multiplier-boundedness}
gives $\|\blmbd^j\|,\|\bmu^j\|,\|\blmbd_j^\star(\bz)\|\le M$ and
$\|\bmu^{k+1}-\bmu^k\|\le2M\beta_k$. Direct comparison of the two
quadratic damping terms gives, uniformly in $\bx$ and $\bz$,
\[
|d_{k+1}(\blmbd;\bz)-d_k(\blmbd;\bz)|
\le\Delta_k,
\qquad
|p_{k+1}(\bz)-p_k(\bz)|\le\Delta_k,
\]
where
\[
\Delta_k:=2M^2\left(\frac1{s_k}-\frac1{s_{k+1}}\right)
+\frac{2M}{s_k}\|\bmu^{k+1}-\bmu^k\|.
\]
Consequently,
\[
|R_k|\le4\Delta_k
\le8M^2\left(\frac1{s_k}-\frac1{s_{k+1}}\right)
+16M^2\frac{\beta_k}{s_k}.
\]
The two explicit positive terms above satisfy
\[
\frac{\beta_k(2-\beta_k)}{2s_k}
\|\blmbd^{k+1}-\bmu^k\|^2
\le4M^2\frac{\beta_k}{s_k},
\]
and
\[
\frac12\left(\frac1{s_k}-\frac1{s_{k+1}}\right)
\|\blmbd^{k+1}-\bmu^{k+1}\|^2
\le2M^2\left(\frac1{s_k}-\frac1{s_{k+1}}\right).
\]
Adding the estimates proves \eqref{eq:full-lyapunov-descent}.

For the stated schedule, the two error series satisfy
\[
\sum_{k=0}^\infty\left(\frac1{s_k}-\frac1{s_{k+1}}\right)=\frac1{c_s},
\qquad
\sum_{k=0}^\infty\frac{\beta_k}{s_k}
\le\frac{4\theta_\mu}{c_s},
\]
which makes the accumulated transport error finite. The lower bound
\eqref{eq:current-lyapunov-lower-bound} follows as follows. For any
$\blmbd\in\partial h(\calA(\bx))$, proximal optimality gives
$\by(\bx,\blmbd)=\calA(\bx)$ and hence
$\mathcal D_\rho(\bx,\blmbd)=F(\bx)$. Since
$\|\blmbd\|\le L_h$ and $\|\bmu^k\|\le M$,
\[
\max_{\blmbd}\mathcal P_k(\bx,\blmbd;\bz)
\ge F(\bx)+\frac1{2\gamma}\|\bx-\bz\|^2
-\frac{(L_h+M)^2}{2s_k}.
\]
Minimizing in $\bx$ and using $s_k\ge s_0$ proves
\eqref{eq:current-lyapunov-lower-bound}, because
$\Psi_k=(\mathcal P_k-d_k)+(p_k-d_k)+p_k\ge p_k$.
Summing \eqref{eq:full-lyapunov-descent} and using the error budget
then gives
\[
\begin{aligned}
&\sum_{k=0}^{N-1}\left[
\frac1{4\eta_k}\|\Delta_{\bx}^{k+1}\|^2
+\frac1{8\tau_k}\|\Delta_{\blmbd}^{k+1}\|^2
+\frac{\alpha_k}{2\gamma}\|\bx^{k+1}-\bz^k\|^2\right] \\
&\quad\le
\Psi_0(\bx^0,\blmbd^0;\bz^0)-\inf_{\bx\in\RR^n}F(\bx)
+\frac{(L_h+M)^2}{2s_0}
+\frac{(10+80\theta_\mu)M^2}{c_s},
\qquad N\ge1.
\end{aligned}
\]
\end{proof}



\begin{proof}[Proof of Theorem~\ref{thm:main}]
The schedule \eqref{eq:ddlalm-parameter-schedule} satisfies
\eqref{eq:parameter-choice}: $\eta_kL_x=1/2$,
$\tau_kL_{d,k}\le\tau_kL_{d,0}\le1/2$,
$\tau_kG_{\cal A}^2\le\mu_\gamma^2/(72L_x)
=\eta_k\mu_\gamma^2/36$, and the two bounds on $\alpha_k$ follow
from $s_k\ge s_0$ and its definition. Moreover,
$s_{k+1}\ge s_k>0$, $\eta_k,\tau_k>0$, and
$0<\beta_k\le1$. Since $L_{d,k}\ge1/\rho+1/s_k$,
the multiplier-boundedness condition also holds.
By assumption, $\inf_{\bx}F(\bx)>-\infty$. The master budget proved after
Lemma~\ref{lem:moving-three-value-descent} is therefore finite and
independent of $K$.

The tail has $\lfloor K/2\rfloor\ge K/3$ elements. Since the primal
and dual stepsizes are fixed and
$\alpha_k\ge\alpha_{K-1}=\Theta(K^{-1/3})$ on the tail, the master
budget and the definition of $\kappa$ give
\[
\|\Delta_{\bx}^{\kappa+1}\|^2=O(K^{-1}),\qquad
\|\Delta_{\blmbd}^{\kappa+1}\|^2=O(K^{-1}),\qquad
\|\bx^{\kappa+1}-\bz^\kappa\|^2=O(K^{-2/3}).
\]
The output certificate follows directly from
\eqref{eq:prox-multiplier}. For stationarity, the $\bx$-update
and $L_{\calD,\bx}$-smoothness of $\mathcal D_\rho(\cdot,\blmbd^k)$ give
\[
\|\nabla f(\bx^{k+1})+J_{\cal A}(\bx^{k+1})^\top\bq^k\|
\le\left(\frac1{\eta_k}+L_x\right)\|\Delta_{\bx}^{k+1}\|
+\frac1\gamma\|\bx^{k+1}-\bz^k\|.
\]
For feasibility, the dual update gives
\[
{\cal A}(\bx^{k+1})-\by^{k+1}
=\frac1{\tau_k}\Delta_{\blmbd}^{k+1}
+\frac1{s_k}(\blmbd^k-\bmu^k),
\]
and hence
\[
\|{\cal A}(\bx^{k+1})-\by^{k+1}\|
\le\frac1{\tau_k}\|\Delta_{\blmbd}^{k+1}\|+\frac{2M}{s_k}.
\]
Square the last two bounds using $(a+b)^2\le2a^2+2b^2$ and use
$s_\kappa^{-2}=O(K^{-2/3})$. Both squared residuals are
$O(K^{-2/3})$, so their norms are $O(K^{-1/3})$. All constants depend
only on the problem data and fixed initialization and parameters, so
$K=O(\epsilon^{-3})$ suffices.
\end{proof}
